% ============================================================
\chapter{Prologue: A Linear Game of Life}
\label{ch:prologue}
% ============================================================

\section*{Chapter Preview}
\addcontentsline{toc}{section}{Chapter Preview}

This prologue gives a preview of the course. It is not meant to be fully
understood on first reading. Its purpose is to show the kind of questions
linear algebra can answer: how a rule moves information, which patterns
persist, and why changing coordinates can make a complicated system simple.

\section{States and Rules}

Imagine a finite grid. At each grid point we store a number. The whole grid is
a state. If the grid has $N$ cells, then a state is a vector in $\R^N$.

A rule takes one state and produces the next state. Some rules are nonlinear,
such as Conway's Game of Life. In this book we focus on linear rules: rules
where
\[
  T(u+v)=T(u)+T(v),
  \qquad
  T(\lambda u)=\lambda T(u).
\]

The linearity condition looks restrictive, but it buys us a great deal. A
linear rule can be studied by understanding what it does to a basis. A linear
rule can be represented by a matrix. Sometimes a linear rule becomes diagonal
after the right change of basis.

\section{A One-Dimensional Toy Rule}

Take a vector $u=(u_1,\ldots,u_N)$ and define a new vector by
\[
  (Tu)_j = \frac{1}{2}u_{j-1}+\frac{1}{2}u_{j+1},
\]
with indices wrapped around the ends. Each entry is replaced by the average of
its two neighbors. This is a linear transformation from $\R^N$ to $\R^N$.

Some patterns are quickly smoothed out. Other patterns keep their shape and
are merely scaled. Those special patterns are eigenvectors. The scaling factors
are eigenvalues. Much later, Fourier modes will give a basis of such patterns
for this kind of translation-invariant rule.

\begin{aimlbox}
Many data and signal-processing systems start with this same idea: a state is
a vector, a rule is a transformation, and the useful coordinates are the ones
where the rule becomes simple. Convolutional filters, smoothing, denoising, and
frequency methods all live near this example.
\end{aimlbox}

\section{The Course Promise}

By the end of the book, the statements above will no longer be mysterious. We
will know what vector spaces are, what linear maps are, why matrices represent
them, how bases change descriptions, and why eigenvectors and singular vectors
are the preferred coordinates for many computations.

\section*{Chapter Summary}
\addcontentsline{toc}{section}{Chapter Summary}

\begin{itemize}
  \item A state with finitely many numbers can be treated as a vector.
  \item A linear rule preserves addition and scalar multiplication.
  \item Eigenvectors are patterns whose direction is preserved by a rule.
  \item A good basis can turn a complicated rule into a simple one.
\end{itemize}

\section*{Where This Appears Later}
\addcontentsline{toc}{section}{Where This Appears Later}

Linear maps begin in Chapter~\ref{ch:matrices-linear-maps}. Eigenvectors and
dynamics appear in Chapter~\ref{ch:eigenvalues-dynamics}. Convolution and the
discrete Fourier transform return in Chapter~\ref{ch:images-convolution-dft}.

\section*{Exercises}
\addcontentsline{toc}{section}{Exercises}

\subsection*{A. Theory}

\begin{exercise}
\exerciseid{1.A1}
Explain why the averaging rule in the chapter is linear. Your answer should
check both addition and scalar multiplication.
\end{exercise}

\begin{exercise}
\exerciseid{1.A2}
Give an example of a rule on grid states that is not linear. State which part
of linearity fails.
\end{exercise}

\subsection*{B. By Hand}

\begin{exercise}
\exerciseid{1.B1}
Let $u=(1,0,0,0)$ on a four-cell circular grid. Apply the averaging rule
$(Tu)_j=(u_{j-1}+u_{j+1})/2$ once by hand.
\end{exercise}

\subsection*{C. Python / NumPy}

\begin{exercise}
\exerciseid{1.C1}
Use basic \Python\ and \NumPy\ to implement the averaging rule for vectors of
length $N$. Test that your function satisfies
\texttt{T(u + v) == T(u) + T(v)} up to numerical tolerance.
\end{exercise}

\subsection*{D. Challenge}

\begin{exercise}
\exerciseid{1.D1}
Find a nonzero vector on a four-cell circular grid that is unchanged by the
averaging rule.
\end{exercise}

